\section{Introduction}
\label{sec:intro}
Gold-standard labels are scarce in many sciences, while features and a trained predictor are cheap. Prediction-powered inference (PPI)~\citep{angelopoulos2023ppi} and its power-tuned version PPI++~\citep{angelopoulos2023ppipp} combine a small labeled sample with predictions on a large unlabeled sample to give confidence intervals that stay valid whatever the quality of the predictor and that are narrower than labeled-only intervals when the predictor is informative.

We ask a question about the \emph{ceiling} of this recipe. Fix the labeled sample $(Y_i,f(X_i))_{i\le n}$, the unlabeled sample $(f(X_j))_{j\le N}$ and the predictor $f$; buy no labels, retrain nothing, add no unlabeled data. \emph{Can a valid 95\% interval for the population mean still be made shorter than the PPI++ interval, by how much, and under which conditions?} We separate three different sources of possible shortening, because each needs its own validity check: (a)~a lower-variance point estimator (efficiency), (b)~a better finite-sample critical value or variance estimate, and (c)~less conservative constructions than CLT-based ones.

Our main object is (a). Within the information set above, the estimand has a semiparametric efficiency bound that depends on the data only through $E[Y\mid f]$ (Section~\ref{sec:setting}); PPI++ attains it exactly when $E[Y\mid f]$ is affine in $f$. We therefore study estimators that learn a calibration map $g\approx E[Y\mid f]$ from the labels with cross-fitting and plug it into the PPI++ construction, and we measure on real and synthetic data, with known truth and thousands of replicates per cell, how much width this buys and whether coverage survives.

\paragraph{Answer in brief.} (1)~On the nine real datasets shipped with \texttt{ppi-python} the answer is essentially \emph{no}: $E[Y\mid f]$ is nearly affine in the predictor, the efficiency ceiling lies within $0.5\%$ of the PPI++ width on seven sets (for gene expression and census income it is $0.7\%$ and $0.2\%$ below the width of PPI++ without its clip), and a cross-fitted calibrated-score interval is wider than PPI++ in 24 of 26 primary cells, by up to $11\%$ at $n=50$, because learning the calibration map from the labels costs more than it gains. (2)~Where $E[Y\mid f]$ is clearly nonlinear (synthetic populations) the calibrated interval can be shorter, by up to $21\%$ in our grids, at coverage no better than PPI++'s in that cell (never below the ceiling), with the gain growing with $n$, $N/n$ and signal-to-noise. (3)~At small $n$ the reference PPI++ interval itself undercovers (down to $0.84$ on census income at $n=50$); a jackknife-$t$ correction restores coverage on most sets for $10$--$15\%$ more width, so ``shorter'' claims at small $n$ must be read against a reference that is not valid. (4)~Under label--unlabeled distribution shift the unclipped power-tuned weight makes every calibrated or unclipped interval collapse (coverage $0.49$ at $n=400$ and $0.18$ at $n=1000$ on plankton) while the clipped package default survives.

\paragraph{Contributions.} We (i)~separate three sources of possible shortening and state the efficiency ceiling for the fixed-input setting with its assumptions (Proposition~\ref{prop:floor}); (ii)~measure it with $2000$ paired replicates per cell on nine real datasets and a one-factor-at-a-time synthetic grid with known truth, reporting coverage with Wilson intervals and paired width ratios with bootstrap intervals; (iii)~show that the measured gap to the ceiling is the plug-in cost of learning the calibration map and quantify it; (iv)~document small-$n$ undercoverage of the plug-in PPI++ interval and the width price of repairing it, and the shift failure of the unclipped weight. We do not propose a new estimator class: the calibrated estimator is the standard augmented estimator with a cross-fitted one-dimensional nuisance.
